\cleardoublepage
\chapter{CONTOH ARTEFAK DATA INSPEKSI KONTAINER}

Lampiran ini menyajikan contoh artefak data yang digunakan untuk memperjelas
struktur informasi pada rancangan dan realisasi sistem inspeksi kontainer.
Contoh yang ditampilkan bukan dokumentasi seluruh kode program, melainkan
cuplikan representatif dari data yang dikirimkan komponen \textit{edge} ke
layanan aplikasi agar hasil identifikasi dan hasil inspeksi dapat dikaitkan
dengan entitas kontainer yang sama.

\section{Contoh \textit{Payload} Identifikasi Kontainer dari \textit{Edge} ke Layanan API}

\begin{lstlisting}[language={}, numbers=none, basicstyle=\ttfamily\scriptsize, breaklines=true, caption={Contoh payload identifikasi kontainer dari komponen edge ke layanan API}, label=lst:container-identification-payload]
{
  "edgeDeviceId": "edge-gate-01",
  "cameraId": "gate-camera-01",
  "capturedAt": "2026-05-20T09:15:30Z",
  "containerNumber": "MSCU1234567",
  "normalizedContainerNumber": "MSCU1234567",
  "ocr": {
    "confidence": 0.92,
    "sourceFrame": "frames/MSCU1234567-001.jpg"
  },
  "evidence": {
    "imageUrl": "evidence/containers/MSCU1234567/identify.jpg"
  }
}
\end{lstlisting}

Listing~\ref{lst:container-identification-payload} menunjukkan contoh data
identifikasi yang menjadi dasar pembentukan atau pemutakhiran entitas
kontainer. Nomor kontainer yang telah dinormalisasi digunakan sebagai pengikat
utama agar data inspeksi berikutnya tidak berdiri sebagai catatan terpisah.

\section{Contoh \textit{Payload} Hasil Inspeksi Kerusakan}

\begin{lstlisting}[language={}, numbers=none, basicstyle=\ttfamily\scriptsize, breaklines=true, caption={Contoh payload hasil inspeksi kerusakan kontainer}, label=lst:defect-scan-payload]
{
  "containerNumber": "MSCU1234567",
  "scanType": "defect",
  "scannedAt": "2026-05-20T09:16:10Z",
  "findings": [
    {
      "category": "dent",
      "severity": "medium",
      "confidence": 0.87,
      "location": "right_panel",
      "evidenceImage": "evidence/containers/MSCU1234567/defect-001.jpg"
    }
  ],
  "summary": {
    "totalFindings": 1,
    "requiresManualReview": true
  }
}
\end{lstlisting}

Listing~\ref{lst:defect-scan-payload} menunjukkan contoh hasil inspeksi
kerusakan yang dikaitkan kembali ke kontainer yang sama melalui nomor kontainer
terkonfirmasi. Pola ini mendukung prinsip riwayat kontainer yang dibahas pada
Bab IV dan Bab V, yaitu pemisahan antara data terstruktur dan artefak visual
tanpa memutus hubungan keduanya.
